\documentclass[11pt]{article}

\usepackage[margin=0.75in]{geometry}
\usepackage{amsmath}
\usepackage{booktabs}
\usepackage{graphicx}
\usepackage{multirow}
\usepackage{tabularx}

\title{AppWorld Topic Distributions and Scenario Samples}
\author{}
\date{}

\begin{document}
\maketitle

\section{Train and test topic distributions}

An AppWorld scenario is a family of three task variants sharing an ID prefix.
The full train, normal-test, and challenge-test splits contain 30, 56, and 139
scenarios, respectively. We describe task content using ten instruction-derived
topics: Amazon, Gmail, Todoist, Splitwise, Spotify, Venmo, Simple Note, Phone,
File system, and Other. App topics require an explicit app mention; Phone and
File system also use corresponding instruction references. An instruction
without any of these tags is assigned Other. This is a coarse description of
task content: for example, an email task that does not explicitly mention
Gmail may be tagged Other.

Each scenario has total weight one. Its three variants each contribute one
third of that weight, divided equally among all topics tagged in that
variant. Let \(T_{s,v}\) be the set of tags in variant \(v\) of scenario \(s\).
For a split \(G\), the topic probability is
\begin{equation}
  p_{G,k}
  = \frac{1}{|G|}
    \sum_{s\in G}\frac{1}{3}
    \sum_{v=1}^{3}
    \frac{\mathbf{1}\{k\in T_{s,v}\}}{|T_{s,v}|},
  \qquad
  \sum_{k=1}^{10}p_{G,k}=1.
  \label{eq:topic-distribution}
\end{equation}

The indicator \(\mathbf{1}\{k\in T_{s,v}\}\) is one if topic \(k\) is
tagged in that variant and zero otherwise. The denominator
\(|T_{s,v}|\) is the number of tags in the variant; it is at least one
because an untagged instruction receives Other. For example, if
\(T_{s,v}=\{\text{Spotify}\}\), the fraction is one for Spotify and zero
for every other topic. If \(T_{s,v}=\{\text{Spotify},\text{Phone}\}\), it is
\(1/2\) for each of those two topics and zero for the rest. That variant
therefore contributes \((1/3)(1/2)=1/6\) of its scenario's weight to each
topic, before averaging over the \(|G|\) scenarios.

\begin{table}[htbp]
\centering
\small
\setlength{\tabcolsep}{4pt}
\caption{Topic probabilities in the full splits, quick2h samples (five
scenarios per split), and intermediate samples (10 train, 19 normal, and 46
challenge scenarios). Each column sums to 100\% before rounding.}
\label{tab:topic-probabilities}
\begin{tabular}{lrrr@{\hspace{10pt}}rrr@{\hspace{10pt}}rrr}
\toprule
Topic & \multicolumn{3}{c}{Train} & \multicolumn{3}{c}{Test normal}
      & \multicolumn{3}{c}{Test challenge} \\
\cmidrule(lr){2-4}\cmidrule(lr){5-7}\cmidrule(lr){8-10}
 & Full & Quick2h & Intermediate & Full & Quick2h & Intermediate
 & Full & Quick2h & Intermediate \\
\midrule
Amazon      & 0.00\%  & 0.00\%  & 0.00\%  & 0.00\%  & 0.00\%  & 0.00\%  & 43.65\% & 43.33\% & 44.20\% \\
Gmail       & 0.00\%  & 0.00\%  & 0.00\%  & 0.00\%  & 0.00\%  & 0.00\%  & 11.87\% & 10.00\% & 11.96\% \\
Todoist     & 0.00\%  & 0.00\%  & 0.00\%  & 3.57\%  & 0.00\%  & 2.63\%  & 0.72\%  & 0.00\%  & 0.00\% \\
Splitwise   & 0.00\%  & 0.00\%  & 0.00\%  & 3.57\%  & 10.00\% & 5.26\%  & 1.08\%  & 0.00\%  & 1.09\% \\
Spotify     & 50.00\% & 50.00\% & 50.00\% & 21.43\% & 20.00\% & 21.05\% & 6.35\%  & 6.67\%  & 5.80\% \\
Venmo       & 20.00\% & 20.00\% & 20.00\% & 34.52\% & 30.00\% & 34.21\% & 7.19\%  & 10.00\% & 7.61\% \\
Simple Note & 8.33\%  & 10.00\% & 5.00\%  & 18.75\% & 20.00\% & 18.42\% & 3.24\%  & 0.00\%  & 3.26\% \\
Phone       & 8.33\%  & 0.00\%  & 10.00\% & 9.23\%  & 10.00\% & 10.53\% & 6.83\%  & 0.00\%  & 6.52\% \\
File system & 13.33\% & 20.00\% & 15.00\% & 7.14\%  & 10.00\% & 7.89\%  & 10.43\% & 10.00\% & 10.87\% \\
Other       & 0.00\%  & 0.00\%  & 0.00\%  & 1.79\%  & 0.00\%  & 0.00\%  & 8.63\%  & 20.00\% & 8.70\% \\
\bottomrule
\end{tabular}
\end{table}

In the full splits, normal test has relatively more Venmo and Simple Note
content and introduces Todoist and Splitwise topics. Challenge test has
substantial Amazon and Gmail content, both absent from the training
instructions under these tags.

\section{Total variation distance}

For two topic distributions \(P=(p_1,\ldots,p_{10})\) and
\(Q=(q_1,\ldots,q_{10})\), total variation (TV) distance is
\begin{equation}
  \operatorname{TV}(P,Q)
  = \frac{1}{2}\sum_{k=1}^{10}|p_k-q_k|.
  \label{eq:tv}
\end{equation}
The factor \(1/2\) keeps the result between zero and one. A value of
\(0\) means the two distributions assign exactly the same probability to
every \emph{measured topic}; their individual tasks need not be identical.
A value of \(1\) means the distributions have disjoint topic support:
no measured topic has positive probability in both. Equivalently, TV is
the minimum fraction of topic probability mass that must be reassigned to
make the distributions match.

Applying Eq.~\eqref{eq:tv} to the full-split proportions in
Table~\ref{tab:topic-probabilities} gives
\begin{align}
  \operatorname{TV}(P_{\mathrm{train}},P_{\mathrm{normal}})
    &= \tfrac{1}{2}(0.695238)=\mathbf{0.347619}, \\
  \operatorname{TV}(P_{\mathrm{train}},P_{\mathrm{challenge}})
    &= \tfrac{1}{2}(1.318945)=\mathbf{0.659472}.
\end{align}
These calculations use the unrounded probabilities; the displayed topic
percentages are rounded to two decimal places. The challenge split is
therefore much farther from train on this topic representation than normal
test is.

\section{Two-hour sampling methods}

\paragraph{Quick2h sample.}
From each source split, select five complete scenarios and retain all three
variants of each scenario: 15 scenarios and 45 tasks overall. Selection
balances the split's instruction-derived topic mix, task difficulty, and number
of apps. The normal and challenge samples also span different levels of
instruction similarity to training scenarios and include test content absent
from train. At the pilot estimate of 2.2 minutes per task, 45 tasks take about
99 minutes.

\paragraph{Conditioned sample.}
Select five complete scenarios per split under an additional memory-pressure
eligibility requirement. Each family is scored from its task variants and
evaluator conditions on memory structure (M), retrieval (I), adaptation (A),
evidence sampling (P), and memory update (U), with scores in
\(\{0,1,2\}\). A family is eligible only if all five scores are at least
one. Among reviewed eligible families, selection aims to retain each original
split's topic mix, difficulty, app count, and test novelty mix. Baseline
success or failure is not used as a screening criterion. The resulting
15 scenarios again contain 45 tasks. Conditioning and the five-scenario
budget can change the observed topic proportions, so neither sample is an
unbiased estimate of full-split performance.

\subsection{Reusable total-variation-preserving selection criteria}

For a new sample, fix the source splits and topic rule, then select complete
three-task families within the budget. The ordinary pool is each full split;
the conditioned pool contains families passing all five tests in
Table~\ref{tab:memory-pressure-tests}. In both cases, compare each sample
against its \emph{full} source split for topic, difficulty, app-count, and
test-novelty-bin TV. Also compare each sampled train-to-test topic TV with its
full-split counterpart: for \(t\in\{\mathrm{normal},\mathrm{challenge}\}\),
\(\delta_t=|\operatorname{TV}(P_{S_{\mathrm{train}}},P_{S_t})-
\operatorname{TV}(P_{F_{\mathrm{train}}},P_{F_t})|\), where \(F\) denotes a
full split and \(S\) its sample. Prefer small \(\delta_t\), retain the
challenge-greater-than-normal ordering and test-only app content where
feasible, and report any mismatch after conditioning. Detailed criteria and
audit fields are in \texttt{samples/memory\_pressure\_screening.md}.

\clearpage
\subsection{Memory-pressure scoring tests}

Table~\ref{tab:memory-pressure-tests} converts the reusable conditioned-sample
screening rubric into score definitions. M, I, A, P, and U denote memory
structure, retrieval, adaptation, evidence sampling, and memory update.
Inspect all three variants and intrinsic evaluator conditions together. Each
dimension receives exactly one score from \(\{0,1,2\}\), supported by a
concrete family observation and a one-sentence justification.

\begin{table}[htbp]
\centering
\footnotesize
\setlength{\tabcolsep}{5pt}
\renewcommand{\arraystretch}{1.22}
\caption{Memory-pressure scoring tests used by the conditioned sample.}
\label{tab:memory-pressure-tests}
\begin{tabularx}{\textwidth}{@{}c
  >{\raggedright\arraybackslash}X
  >{\raggedright\arraybackslash}X
  >{\raggedright\arraybackslash}X@{}}
\toprule
Mode & Score = 0 & Score = 1 & Score = 2 \\
\midrule
\textbf{M}
  & No meaningful reusable procedure, or variants are essentially identical.
  & Variants share a procedure but differ in a task-specific binding, entity, condition, threshold, or object set.
  & Score 1 holds, and transferring a variant-specific detail could cause an error; reusable invariants must be separated from scoped details. \\
\addlinespace
\textbf{I}
  & Only one obvious reusable procedure is needed, with no meaningful competing memories.
  & Multiple relevant reusable rules are needed, or plausible but inapplicable distractor rules exist.
  & Multiple relevant rules must be retrieved jointly and plausible distractors exist; missing or misretrieving one changes the procedure. \\
\addlinespace
\textbf{A}
  & A retrieved generic rule can be executed almost directly.
  & At least one non-trivial task-specific binding, filter, entity resolution, or ordered action is required.
  & At least two non-trivial operations must be composed, such as predicates, set operations, aggregation, dependencies, an output contract, or a postcondition check. \\
\addlinespace
\textbf{P}
  & Failure feedback is binary or ambiguous and cannot provide a useful diagnosis.
  & Evaluator or trajectory evidence localizes a failure and supports a reusable lesson.
  & Feedback identifies a root cause amid plausible competing trajectory explanations; learning requires selecting the causal evidence. \\
\addlinespace
\textbf{U}
  & Earlier variants offer little or no reusable information for later variants.
  & At least one earlier lesson should transfer to a later variant.
  & Both transferable and variant-specific non-transferable information exist, testing useful updates and negative transfer. \\
\bottomrule
\end{tabularx}
\end{table}

A family is a pressure-test candidate only when all five scores are at least
one. Do not award a positive score from task difficulty alone or use baseline
success and failure. For P, use feedback available to the intended learner;
score 2 requires an actual trajectory with plausible competing explanations.

\clearpage
\section{Summary}

\begin{table}[htbp]
\centering
\small
\setlength{\tabcolsep}{5pt}
\caption{Full splits, quick2h and conditioned five-scenario-per-split samples,
and the one-third intermediate sample. Topic TV compares each row to its own
full source split; it is zero for the full splits by definition. The two
rightmost columns compare train and test topic distributions \emph{within the
same set}; each value is shown once per three-row set.}
\label{tab:summary}
\resizebox{\textwidth}{!}{%
\begin{tabular}{llrrrrr}
\toprule
Sets & Splits &
\shortstack{Number of\\Scenario} &
\shortstack{Number of\\Tasks} &
Topic TV &
\shortstack{Train $\rightarrow$\\Normal TV} &
\shortstack{Train $\rightarrow$\\Challenge TV} \\
\midrule
\multirow{3}{*}{Full}
  & Train     & 30  & 90  & 0.0000 & \multirow{3}{*}{0.3476} & \multirow{3}{*}{0.6595} \\
  & Normal    & 56  & 168 & 0.0000 & & \\
  & Challenge & 139 & 417 & 0.0000 & & \\
\midrule
\multirow{3}{*}{Quick2h Sample}
  & Train     & 5 & 15 & 0.0833 & \multirow{3}{*}{0.4000} & \multirow{3}{*}{0.7333} \\
  & Normal    & 5 & 15 & 0.1131 & & \\
  & Challenge & 5 & 15 & 0.1448 & & \\
\midrule
\multirow{3}{*}{Intermediate}
  & Train     & 10 & 30  & 0.0333 & \multirow{3}{*}{0.3605} & \multirow{3}{*}{0.6594} \\
  & Normal    & 19 & 57  & 0.0374 & & \\
  & Challenge & 46 & 138 & 0.0159 & & \\
\midrule
\multirow{3}{*}{Conditioned Sample}
  & Train     & 5 & 15 & 0.1000 & \multirow{3}{*}{0.3000} & \multirow{3}{*}{0.7000} \\
  & Normal    & 5 & 15 & 0.1131 & & \\
  & Challenge & 5 & 15 & 0.2233 & & \\
\bottomrule
\end{tabular}%
}
\end{table}

All three sampled sets preserve the ordering of the full data: challenge is farther
from train than normal test. These distances concern the ten-topic
representation, not the complete distribution of scenario semantics.

\end{document}
